\section{Related work}\label{sec:related}

Valiant introduced the algebraic complexity framework and \(\VNP\)
\cite{Valiant1979,Valiant1980}.  Ben-Or and Cleve showed that polynomial-size formulas and
width-three ABPs are polynomially equivalent \cite{BenOrCleve1992}.  Allender and Wang
proved explicit limitations of exact width-two ABPs \cite{AllenderWang2016}.

Bringmann, Ikenmeyer and Zuiddam studied small-width ABPs, border complexity, and
nondeterministic closures in a unified framework
\cite{BringmannIkenmeyerZuiddam2018,BringmannIkenmeyerZuiddam2017}.  Their Theorem~4.2
identifies \(\VNPone\) with \(\VNP\) when the characteristic is not two; their support-two
refinement appears as Theorem~6.3 in the extended version.  Their Proposition~9.1 proves
strictness over \(\F_2\), while their discussion explicitly leaves all other
characteristic-two fields open.  They also note an \(\F_4\) representation of \(1+xy\),
which shows that the \(\F_2\) obstruction does not automatically extend to larger
characteristic-two fields.  The present \(\tau\)-gadget works uniformly without requiring
an \(\F_4\) subfield.

Dutta et al. proved universality results for \emph{border} width-two ABPs in
characteristic two and showed that a suitable power of a polynomial with a small formula
has a small border-width-two representation \cite{DuttaEtAl2024}.  That work concerns
approximate deterministic width two.  The present theorem concerns exact nondeterministic
width one; neither result subsumes the other.

\section{Limitations and open directions}\label{sec:open}

\paragraph{Constants.}
The certified list keeps a unit placeholder for every ternary vertex and uses ordered pair
exclusions.  Removing redundant unit factors or using a different selector may improve
\(14\) and \(34\).  Such improvements would not change the class theorem.

\paragraph{Local minimality.}
The artifact proves that the characteristic-two gadget achieves one auxiliary bit and three
factors, and certifies a limited degree lower bound in the zero-auxiliary setting.  It does
not formalize the full impossibility of a two-factor positive-auxiliary gadget.  No
optimality claim is part of the main theorem.

\paragraph{The exceptional field.}
Over \(\F_2\), unrestricted nondeterministic width one is strictly weaker than \(\VNP\).
The internal structure of support-bounded subclasses over \(\F_2\) remains a separate
question.

\paragraph{Alternative direct compilers.}
The graph selector gives a modular and machine-checked proof.  A direct activation-state
compiler on the formula tree might reduce constants and shorten the semantic argument, but
such a route should be treated as a new proof until separately formalized and audited.

\paragraph{Border complexity.}
The exact hypercube identities do not provide a deterministic width-two matrix product,
control an approximation parameter, or extract Frobenius roots.  Extending them to the
border-width-two setting would require an additional theorem, not a reinterpretation of
the present proof.

\section{Conclusion}

Over every characteristic-two field other than \(\F_2\), a binary arithmetic formula can
be compiled into a Boolean hypercube sum of a product of affine forms, each supported on at
most two variables, with linear overhead.  The construction works because a uniform
quadratic gadget and a ternary gadget can be integrated with a path selector whose local
exclusions annihilate the ternary error pointwise.  The resulting class equality fills the
characteristic-two range left open in the width-one nondeterministic programme and, together
with prior work, makes \(\F_2\) the unique exceptional field.

\appendix

\section{A more explicit selector proof}\label{app:selector}

For completeness, we isolate the combinatorial fact used in \cref{prop:selector}.  Let
\(H\) be the selected-edge subgraph.  Pair exclusions imply
\(\indeg_H(v),\outdeg_H(v)\in\{0,1\}\).  A nonzero source factor gives
\(\outdeg_H(s)=1\), and a nonzero sink factor gives \(\indeg_H(t)=1\).  At every other
vertex, the flow factor is nonzero if and only if
\(\indeg_H(v)=\outdeg_H(v)\).

Follow the unique selected outgoing edge from \(s\).  At any internal vertex reached, the
selected incoming edge forces a selected outgoing edge.  Because rank strictly increases,
the walk cannot repeat a vertex and must terminate.  It cannot terminate internally and
therefore terminates at \(t\).  Now suppose another selected edge exists.  Its component
contains neither \(s\) nor \(t\), since each of those already has its unique selected
incident edge in the constructed path.  Every vertex in the extra component has equal
selected indegree and outdegree.  Starting from any selected edge and repeatedly following
selected outgoing edges produces a directed cycle in the finite component, contradicting
strict rank increase.  Hence the selected subgraph is exactly one source--sink path.

\section{The exact compiler ledger}\label{app:ledger}

Let \(E,V,P,D\) be as in \cref{sec:compiler}.  The formal compiler uses the following
objects.

\begin{center}
\begin{tabular}{@{}lll@{}}
\toprule
Object & Count & Contribution\\
\midrule
Original edge selectors & \(|E|\) & auxiliary bits\\
Quadratic occurrences & \(|E|+P\) & one bit, three factors each\\
Ternary occurrences & \(D\) & one bit, four factors each\\
Flow placeholders & \(|V|\) & one factor each\\
\bottomrule
\end{tabular}
\end{center}

Therefore
\[
 q=|E|+(|E|+P)+D=2|E|+P+D,
\]
and
\[
 M=|V|+3(|E|+P)+4D.
\]
For formula graphs, substituting
\[
 |E|=l+4a+u,\quad |V|=2l+2a,\quad P=4a,\quad D\le|V|
\]
gives the sharp certified inequalities in \cref{sec:counts}.  The unproved proposed
identity \(D+B=2a\) is not used: the number of degree-three vertices depends on formula
shape beyond the triple \((l,a,u)\).

\section{Machine-checked scope firewalls}\label{app:firewalls}

The formal development includes negative or limiting statements designed to prevent the
main theorem from being overread.

\begin{enumerate}[leftmargin=1.8em]
\item \(abc(1+ab)\) is not the zero formal polynomial.
\item The ternary gadget without a pair exclusion retains a nonzero error.
\item \(\tau=0\) and \(\tau=1\) make \(\kappa=0\), so the normalized gadgets are illegal.
\item Sharing one fresh bit between distinct gadget occurrences changes the represented
polynomial in general.
\item A graph edge carrying an unrestricted three-variable affine label can violate the
support-two target; the theorem is stated for formula-generated graphs with one-variable
or constant labels.
\item A degree-three flow factor genuinely has support three before ternary compilation.
\item The class theorem is a representation theorem and includes no evaluator or runtime
claim.
\end{enumerate}

\begin{thebibliography}{BIZ18}

\bibitem[BIZ18]{BringmannIkenmeyerZuiddam2018}
K.~Bringmann, C.~Ikenmeyer, and J.~Zuiddam.
\newblock On algebraic branching programs of small width.
\newblock \emph{Journal of the ACM}, 65(5):32:1--32:29, 2018.
\newblock \href{https://doi.org/10.1145/3209663}{doi:10.1145/3209663}.

\bibitem[BIZ17]{BringmannIkenmeyerZuiddam2017}
K.~Bringmann, C.~Ikenmeyer, and J.~Zuiddam.
\newblock On algebraic branching programs of small width.
\newblock In \emph{32nd Computational Complexity Conference (CCC 2017)},
  volume 79 of \emph{LIPIcs}, pages 20:1--20:31, 2017.
\newblock \href{https://doi.org/10.4230/LIPIcs.CCC.2017.20}{doi:10.4230/LIPIcs.CCC.2017.20}.

\bibitem[Val79]{Valiant1979}
L.~G. Valiant.
\newblock Completeness classes in algebra.
\newblock In \emph{Proceedings of the Eleventh Annual ACM Symposium on Theory of Computing},
  pages 249--261, 1979.
\newblock \href{https://doi.org/10.1145/800135.804419}{doi:10.1145/800135.804419}.

\bibitem[Val80]{Valiant1980}
L.~G. Valiant.
\newblock Reducibility by algebraic projections.
\newblock Internal report, Department of Computer Science, University of Edinburgh, 1980.

\bibitem[BC92]{BenOrCleve1992}
M.~Ben-Or and R.~Cleve.
\newblock Computing algebraic formulas using a constant number of registers.
\newblock \emph{SIAM Journal on Computing}, 21(1):54--58, 1992.
\newblock \href{https://doi.org/10.1137/0221006}{doi:10.1137/0221006}.

\bibitem[AW16]{AllenderWang2016}
E.~Allender and F.~Wang.
\newblock On the power of algebraic branching programs of width two.
\newblock \emph{Computational Complexity}, 25(1):217--253, 2016.
\newblock \href{https://doi.org/10.1007/s00037-015-0114-7}{doi:10.1007/s00037-015-0114-7}.

\bibitem[DIK${}^{+}$24]{DuttaEtAl2024}
P.~Dutta, C.~Ikenmeyer, B.~Komarath, H.~Mittal, S.~G. Nanoti, and D.~Thakkar.
\newblock On the power of border width-2 {ABP}s over fields of characteristic 2.
\newblock In \emph{41st International Symposium on Theoretical Aspects of Computer Science
  (STACS 2024)}, volume 289 of \emph{LIPIcs}, pages 31:1--31:16, 2024.
\newblock \href{https://doi.org/10.4230/LIPIcs.STACS.2024.31}{doi:10.4230/LIPIcs.STACS.2024.31}.

\bibitem[Bür00]{Burgisser2000}
P.~Bürgisser.
\newblock \emph{Completeness and Reduction in Algebraic Complexity Theory}.
\newblock Algorithms and Computation in Mathematics, vol.~7. Springer, 2000.
\newblock \href{https://doi.org/10.1007/978-3-662-04179-6}{doi:10.1007/978-3-662-04179-6}.

\bibitem[MP08]{MalodPortier2008}
G.~Malod and N.~Portier.
\newblock Characterizing Valiant's algebraic complexity classes.
\newblock \emph{Journal of Complexity}, 24(1):16--38, 2008.
\newblock \href{https://doi.org/10.1016/j.jco.2006.09.006}{doi:10.1016/j.jco.2006.09.006}.

\bibitem[dMU21]{Lean4}
L.~de Moura and S.~Ullrich.
\newblock The {Lean 4} theorem prover and programming language.
\newblock In \emph{Automated Deduction---CADE 28}, volume 12699 of \emph{LNCS},
  pages 625--635. Springer, 2021.
\newblock \href{https://doi.org/10.1007/978-3-030-79876-5_37}{doi:10.1007/978-3-030-79876-5\_37}.

\bibitem[Com20]{Mathlib2020}
The mathlib Community.
\newblock The {Lean} mathematical library.
\newblock In \emph{Proceedings of the 9th ACM SIGPLAN International Conference on
  Certified Programs and Proofs}, pages 367--381, 2020.
\newblock \href{https://doi.org/10.1145/3372885.3373824}{doi:10.1145/3372885.3373824}.

\end{thebibliography}
